% ── Formal Echo: Bayes ────────────────────────────────────────
% Appears after Mackandal's distributed network.
% Tone: provisional, iterative, always in motion, never arriving.

\begin{formalecho}{Thomas Bayes}{The Infinite Approach}

He is not certain.

This is not a failure. He has come to understand it as something
closer to a starting condition---the only honest description of
where any process of reasoning must begin.

He has a belief. It is not arbitrary; it is founded on what he has
seen, what he has been told, what has held up under prior
examination. But it is not fixed. It is a weight assigned to a
possibility, and weights can be revised.

This is the part that others find difficult.

They want the process to terminate. They want a moment at which
the reasoning is complete, the conclusion declared, the question
closed. He understands the desire. He has felt it himself. But he
has also noticed what happens when closure is forced before the
evidence permits it: the system becomes brittle, resistant to what
comes next, committed to a position that may need to be abandoned.

Better to hold the belief loosely.

Better to say: given what I have seen, this is how likely it seems.
And when I see more, I will adjust.

\medskip

The adjustment is the point.

Not the belief at any given moment, but the process by which it
moves. Each new observation enters the system and shifts the
weights---slightly, sometimes, almost imperceptibly; dramatically,
occasionally, when what arrives is sufficiently unexpected.

The belief approaches something.

He is careful about how he names that something. Not truth,
exactly---that word carries too much finality. Something more like:
the configuration that the evidence, accumulated, is pressing
toward. A limit that the process asymptotically nears without
necessarily reaching.

He finds this neither troubling nor consoling.

It is simply the shape of the thing.

\medskip

What he cannot fully resolve is the question of the prior.

Every process must begin somewhere. Every assignment of initial
weights reflects something---experience, assumption, the particular
angle from which the problem was first approached. Two people
beginning from different priors, receiving identical evidence, will
arrive at different positions. They will converge, given enough
evidence. But convergence takes time, and time is not always
available.

He turns this over without reaching a conclusion.

It is possible that the prior cannot be eliminated---that every
system of reasoning carries within it the residue of where it
began, irreducibly, however much evidence is accumulated.

If so, then no process of updating, however rigorous, achieves
full independence from its starting point.

The belief moves. It approaches. It does not fully arrive.

He writes this down, then sets the pen aside.

Outside, something is happening that he cannot see from where he
sits. The evidence for it has not yet reached him. When it does,
he will adjust.

Until then, he holds what he has.

Loosely.

With appropriate weight.

Ready to revise.

\end{formalecho}
